Для повышения качества последующего распознавания в разработанном программном обеспечении реализован этап адаптивной предобработки изображений страниц. Его назначение состоит в том, чтобы до запуска модулей распознавания уменьшить влияние типичных дефектов сканирования выделенным Akhil~\cite{refAkhil}. В отличие от фиксированной схемы фильтрации, в разработанном алгоритме параметры обработки не задаются вручную заранее, а вычисляются автоматически на основе характеристик конкретного входного изображения.

Работа алгоритма начинается с загрузки изображения страницы и его преобразования в градации серого. Далее для получения грубой бинарной маски используется пороговая обработка по методу Отсу с инверсией яркости описанный описанного в работе Старовойтова~\cite{refStarovoytov}. Полученное бинарное изображение применяется не как финальный результат, а как промежуточный этап для оценки характерного размера текстовых элементов. На этой стадии выделяются внешние контуры связных компонент, для каждой из которых вычисляется ограничивающий прямоугольник. Из дальнейшего анализа исключаются слишком малые компоненты, соответствующие шуму, а также крупные объекты, которые, как правило, относятся к линиям, рамкам или фрагментам графики. После фильтрации по высоте и ширине для оставшихся компонент вычисляется медианная высота, интерпретируемая как оценка характерной высоты символа на странице.

Полученная оценка используется в качестве основного управляющего параметра всей последующей предобработки. Через оценку для размера одного сиволап изображения масштабируется до 40 пикселей --- эмпирически выбранная целевая высота символа, обеспечивающая устойчивую работу современных OCR-движков также в программе реализовано предотвращение чрезмерного увеличения изображения и роста вычислительных затрат. Таким образом, страницы с мелким текстом автоматически масштабируются сильнее, а страницы, уже имеющие достаточное разрешение, обрабатываются без лишнего увеличения.

На основе той же оценки высоты одного символа адаптивно вычисляются параметры локальной бинаризации. Размер окна бинаризации выбирается пропорционально примерно четырём высотам символа, после чего корректируется до ближайшего нечётного значения и дополнительно ограничивается размерами изображения. Такой выбор обусловлен тем, что окно должно быть достаточно большим для оценки локального фона, но не настолько большим, чтобы сглаживать полезные детали символов и тонких графических элементов. Константа адаптивной бинаризации также вычисляется автоматически как доля от размера окна, что позволяет согласованно изменять чувствительность метода в зависимости от масштаба текста на странице.

Кроме бинаризации, алгоритм автоматически подбирает интенсивность ряда дополнительных операций предобработки. Уровень шумоподавления, усиления контраста и резкости зависит от рассчитанного коэффициента масштабирования: чем меньше исходный размер символов и чем сильнее требуется увеличение изображения, тем более выраженной становится последующая коррекция. Это позволяет избежать как недостаточной обработки слабоконтрастных изображений, так и переусиления качественных сканов. Интенсивность морфологических операций также задаётся пропорционально оценённой высоте символа. Одна группа параметров ориентирована на удаление мелких артефактов и разрывов, другая --- на объединение близко расположенных фрагментов символов, что особенно важно при работе с деградировавшими сканами.

В реализованной схеме по умолчанию активируется коррекция небольшого поворота страницы, поскольку даже малый наклон существенно ухудшает как сегментацию строк текста, так и последующее распознавание. Дополнительно рассчитывается величина внутреннего отступа изображения, пропорциональная высоте символа, что позволяет снизить риск потери граничных элементов при дальнейшей обработке. В показанном варианте алгоритма удаление теней и границ страницы отключено, однако архитектура модуля допускает их включение при необходимости для отдельных классов документов.

Предложенный алгоритм представляет собой эвристически адаптивную схему предобработки, в которой центральную роль играет оценка характерного масштаба текста. Такой подход обеспечивает автоматическую настройку параметров под конкретную страницу без необходимости ручного подбора коэффициентов для каждого изображения. Это особенно важно для архивных химических книг, где качество сканов существенно варьируется от страницы к странице. Полученные после предобработки изображения затем передаются специализированным модулям распознавания: текстовые области --- в OCR-подсистему.
%, а графические фрагменты с формулами и реакциями --- в соответствующие средства анализа химических структур. 
Тем самым предобработка выступает базовым этапом, повышающим устойчивость всего конвейера оцифровки.